\section{Set-Like Relations}

In what follows, let $A$ be a class and $R$ be a class relation on $A$.

The property we want is that looking below any one point of $A$ only ever sees a set's worth of $A$.

\begin{boxdefinition}[Set-Like]
    We say $(A, R)$ is \textbf{set-like} iff
    \begin{align*}
        \forall y \in A, \, \setst{x \in A}{x \, R \, y}
    \end{align*}
    is a set.
\end{boxdefinition}

The example we care about most has this property for the most basic of reasons.

\begin{boxexample}[Membership]
    $(A, \in)$ is set-like for every class $A$.
\end{boxexample}

In fact, the fact in the above example is \textbf{equivalent} to the axiom scheme of comprehension (\Cref{ZFC:Compr})! Indeed, if $A = \setst{x}{\varphi\of{x, \bar{z}}}$ and $y \in A$, then
\begin{align*}
    \setst{x \in A}{x \in y} = \setst{x \in y}{\varphi\of{x, \bar{z}}}
\end{align*}
So ``membership is set-like'' is a fancy way of saying ``Comprehension''.

But not all relations are set-like.

\begin{boxnexample}
    Let $A = \OR$. Define $R$ by the following picture:
    \begin{center}
        \begin{tikzpicture}[every node/.style = {font = \small}, dot/.style = {circle, fill, inner sep = 1.5pt}]
            \node[dot, label = right:{$1$}] (1) at (0, 0) {};
            \node[dot, label = right:{$2$}] (2) at (0, 0.8) {};
            \node (d1) at (0, 1.7) {$\vdots$};
            \node[dot, label = right:{$\omega$}] (w) at (0, 2.6) {};
            \node[dot, label = right:{$\omega + 1$}] (w1) at (0, 3.4) {};
            \node (d2) at (0, 4.3) {$\vdots$};
            \node[dot, label = right:{$0$}] (0) at (0, 5.4) {};
            \draw (1) -- (2);
            \draw (2) -- (d1);
            \draw (d1) -- (w);
            \draw (w) -- (w1);
            \draw (w1) -- (d2);
            \draw (d2) -- (0);
        \end{tikzpicture}
    \end{center}
    That is, it looks a bit like $\parenth{\OR \setminus \set{0}, <}$.
\end{boxnexample}

% 9 September

\begin{boxdefinition}[Extensionality]
    $(A, R)$ is \textbf{extensional} iff $\forall y, z \in A$,
    \begin{align*}
        \setst{x \in A}{x \, R \, y} = \setst{x \in A}{x \, R \, z}
        \to y = z
    \end{align*}
\end{boxdefinition}

This is just like the axiom of extensionality, but worded as a property of sets and relations.

Not all relations are extensional.

\begin{boxnexample}[Non-Extensional Relations] \hfill
    If $A = \set{\set{1}, \set{2}}$, then $(A, \in)$ is not extensional. The reason is that
    \begin{align*}
        \setst{x \in A}{x \in \set{1}} = \emptyset = \setst{x \in A}{x \in \set{2}}
    \end{align*}
    but obviously $\set{1} \ne \set{2}$.
\end{boxnexample}

But many are.

\begin{boxexample}[Extensional Relations] \hfill
    \begin{enumerate}
        \item If $M$ is a transitive class, then $(M, \in)$ is extensional. Indeed, suppose $y, z \in M$ and $y \ne z$. Since $M$ is transitive, $y, z \ssq M$, so $y = y \cap M$ and $z = z \cap M$. Thus
            \begin{align*}
                y \cap M = \setst{x \in M}{x \in y} \neq \setst{x \in M}{x \in z} = z \cap M
            \end{align*}

        \item If $(A, R)$ is a class linear ordering (ie, a linear ordering on the class $A$), then $(A, R)$ is extensional.

        For instance, if $A \ssq \OR$ and $R$ is the usual ordering of ordinals (write $R$ as $\in$ or $<$), then $(A, R)$ is extensional.
    \end{enumerate}
\end{boxexample}

We are now ready to state and prove Mostowski Collapse. Note that it will use neither Choice nor Foundation.

\begin{boxtheorem}[Mostowski Collapse Theorem Scheme]
    Let $(A, R)$ be a class structure that's well-founded, set-like and extensional. Define, using the theorem scheme on definitions by recursion, the mapping $G : A \to V$ by
    \begin{align*}
        \forall y \in A, G(y) = \setst{G(x)}{x \in A \text{ and } x \, R \, y}
    \end{align*}
    Write the range of $G$ as $M$. Then,
    \begin{enumerate}
        \item $M$ is a transitive class
        \item $G$ is a class isomorphism from $(A, R)$ to $(M, \in)$
        \item If $M'$ is a transitive class and $G' : (A, R) \simeq (M, \in)$, then $M' = M$ and $G' = G$.
    \end{enumerate}
    Note that $G$ is referred to as the \textbf{Mostowski Collapse Isomorphism}.
\end{boxtheorem}
\begin{remark}
    When we say that $M' = M$ and $G' = G$ above, we are not saying that their defining formulae are the same. In other words, their equality is not definitional.
\end{remark}
\begin{proof}
    We may assume $R \ssq A \times A$.

    \begin{enumerate}
        \item Say $a \in b \in M$. We must see $a \in M$. Then, $b \in M$. As $M$ is the range of $G$, there's some $y \in  A$ such that $b = G(y)$. As $g(y) = \setst{G(x)}{x \, R \, y}$, there's some $x \, R \, y$ such that $a = G(x)$. Now $a = G(x) \in M$.

        \item We need to show $G$ is an order-preserving bijection.
        
        \begin{itemize}
            \item \underline{Bijectivity.} It's clear that $G$ is surjective, because $M$ is the range of $G$. To show $G$ is injective, it is enough to show that for every $x \in  A$,
            \begin{align}
                \forall x' \in A \brac{G(x) = G\of{x'} \to x = x'}
                \label{Ch2:Eq:Mostowski-pf-expr}
            \end{align}
            Let's prove \eqref{Ch2:Eq:Mostowski-pf-expr} by induction on $R$. Let $y \in A$ and assume \eqref{Ch2:Eq:Mostowski-pf-expr} holds for all $x \, R \, y$. We'll now prove that \eqref{Ch2:Eq:Mostowski-pf-expr} holds for $y$ as well.
    
            Fix any $y' \in A$ and assume that $G(y) = G\of{y'}$. We need to prove $y = y'$. Because $(A, R)$ is extensional, it is enough to show that
            \begin{align*}
                \setst{x}{x \, R \, y} = \setst{x'}{x' \, R \, y}
            \end{align*}
            \begin{description}
                \item[$(\subseteq)$]
                    Suppose $x \, R \, y$. Then
                    \begin{align*}
                        G(x) \in G(y) = G\of{y'} = \setst{G\of{x'}}{x' \, R \, y'}
                    \end{align*}
                    So there's some $x' \, R \, y'$ with $G(x) = G\of{x'}$. By \eqref{Ch2:Eq:Mostowski-pf-expr} for $x$, we know that $x = x'$. Thus $x \, R \, y'$.
    
                \item[$(\supseteq)$]
                    Suppose $x' \, R \, y'$. Then,
                    \begin{align*}
                        G\of{x'} \in G\of{y'} = G(y) = \setst{G(x)}{x \, R \, y}
                    \end{align*}
                    So there's some $x \, R \, y$ such that $G\of{x'} = G(x)$. By \eqref{Ch2:Eq:Mostowski-pf-expr} for $x$, $x' \, R \, y$.
            \end{description}

            \item \underline{Monotonicity.} Suppose $G(x) \in G(y)$. By definition of $G(y)$, there is some $x' \, R \, y$ such that $G(x) = G\of{x'}$. As $G$ is injective, we know that $x = x'$. Thus, $x \, R \, y$.
        \end{itemize}

        \item Suppose $M'$ is a transitive class and $G' : (A, R) \simeq M'$. We need $M = M'$ and $G = G'$.

        Since $M$ is the range of $G$ and $M'$ is the range of $G'$, it is enough to show that $G = G'$. By the uniqueness of recursive definitions, % [CLAUDE] add a cross-reference for this
        it is enough to see that $G'$ satisfies the same defining equation as $G$. That is, $G'(y) = \setst{G'\of{x}}{x \, R \, y}$.

        \begin{description}
            \item[$(\subseteq)$]
                $x \, R \, y \to G'(x) \in G'(y)$ just because $G'$ preserves order (by assumption).

            \item[$(\supseteq)$]
                Suppose $a \in G'(y)$. Now $G'(y) \in M'$ trivially because $M'$ is the range of $G'$. Since $M'$ is a transitive class, $G'(y) \ssq M'$. Hence, $a \in M'$. So we have some $x \in A$ such that $a = G'(x)$. Now $G'(x) \in G'(y)$. Since $G'$ is an isomorphism, this implies $x \, R \, y$.
        \end{description}
    \end{enumerate}
\end{proof}

Let's study some examples.

Suppose $(A, R)$ is a class well-ordering. There are two possibilities:
\begin{enumerate}
    \item $(A, R)$ is a set structure (so a \textit{set} well-ordering).

    If $G : (A, R) \simeq (M, \in)$ is the Mostowski collapse isomorphism,
\end{enumerate}